% Part: lambda-calculus
% Chapter: lambda-definability
% Section: lambda-definable-recursive

\documentclass[../../../include/open-logic-section]{subfiles}

\begin{document}

\olfileid[id]{lam}{dfl}{ldr}
\olsection{Fungsi yang \usetoken{S}{lambda definable} Bersifat Rekursif}

Bukan hanya semua fungsi rekursif parsial bersifat !!{lambda definable};
konversnya juga benar. Artinya, semua fungsi yang !!{lambda definable}
bersifat rekursif parsial.

\begin{thm}
  \ollabel{thm:lambda-computable} Jika suatu fungsi parsial $f$ bersifat
  !!{lambda definable}, maka fungsi itu rekursif parsial.
\end{thm}

\begin{proof}
  Kita hanya memberikan garis besar bukti. Pertama, kita mengaritmetisasi
  suku-$\lambd$, yakni secara sistematis menetapkan bilangan G\"odel bagi
  suku-$\lambd$ dengan menggunakan pengodean barisan melalui perpangkatan
  bilangan prima yang lazim. Kemudian kita mendefinisikan fungsi rekursif
  parsial $\fn{normalize}(t)$ dengan argumen berupa bilangan G\"odel~$t$.
  Dengan pengodean ini, relasi bahwa suatu bilangan mengodekan barisan
  berhingga reduksi satu-langkah yang dimulai dengan suku masukan dan
  berakhir pada suatu bentuk normal bersifat rekursif primitif. Melalui
  minimisasi tak berbatas atas kode barisan semacam itu, fungsi tersebut
  mengembalikan bilangan G\"odel bentuk normal jika bentuk normal itu ada;
  jika tidak, pencariannya tidak berhenti dan fungsi tersebut tidak
  terdefinisi.

  Selanjutnya, definisikan dua fungsi rekursif parsial
  $\fn{toChurch}$ dan $\fn{fromChurch}$ dengan arah yang tegas: fungsi pertama
  memetakan bilangan asli ke bilangan G\"odel numeral Church yang
  bersesuaian, sedangkan fungsi kedua memetakan kode semacam itu kembali ke
  bilangan asli yang direpresentasikannya dan tidak terdefinisi pada kode yang
  tidak mengodekan numeral Church.

  Dengan menggunakan fungsi-fungsi rekursif tersebut, kita dapat
  merepresentasikan fungsi~$f$ sebagai fungsi rekursif parsial. Terdapat suatu
  suku-$\lambd$~$F$ yang !!{lambda define}s~$f$. Untuk menghitung
  $f(n_1, \dots, n_k)$, mula-mula peroleh bilangan G\"odel numeral Church yang
  bersesuaian dengan menggunakan $\fn{toChurch}(n_i)$. Kemudian, berdasarkan
  $\Gn{F}$ dan kode-kode numeral tersebut, gunakan operasi rekursif primitif
  untuk pengodean aplikasi suku guna membentuk bilangan G\"odel dari
  $F \num{n_1}\dots\num{n_k}$. Sekarang terapkan $\fn{normalize}$ pada bilangan
  G\"odel ini. Jika $f(n_1, \dots, n_k)$ terdefinisi,
  $F \num{n_1}\dots\num{n_k}$ mempunyai bentuk normal, yang harus berupa suatu
  numeral Church. Jika tidak, suku tersebut tidak mempunyai bentuk normal,
  sehingga
  \[\fn{normalize}(\Gn{F\num{n_1}\dots\num{n_k}})\]
  tidak terdefinisi. Terakhir, terapkan $\fn{fromChurch}$ pada bilangan
  G\"odel suku yang telah dinormalisasi. Seluruh langkah ini merupakan
  komposisi fungsi rekursif parsial, sehingga fungsi semula rekursif parsial.
\end{proof}

\end{document}
